% =============================================
% Section 1: Introduction
% =============================================
\section{Introduction}
\label{sec:introduction}

The proliferation of large language models (LLMs) has transformed natural
language processing, with models such as GPT-4~\cite{openai2023gpt4},
LLaMA~\cite{touvron2023llama}, and Gemini~\cite{team2023gemini} achieving
human-level fluency across diverse tasks. However, these achievements come at
significant computational cost: even the smallest production-grade models
require gigabytes of memory and substantial floating-point throughput,
effectively restricting deployment to cloud infrastructure and high-end consumer
hardware.

This stands in stark contrast to the billions of deployed microcontrollers
(MCUs) that power IoT devices, wearable electronics, and industrial sensors.
The Espressif \espss{}, a representative modern MCU, features a dual-core
Xtensa LX7 processor at 240\,MHz with only 8\,MB of external \psram{} and up
to 16\,MB of Flash storage. Running any form of coherent language generation on
such hardware was, until recently, considered infeasible.

Recent developments have begun to challenge this assumption. Karpathy's
\texttt{llama2.c}~\cite{karpathy2023llama2c} demonstrated that a minimal
C-based Transformer implementation can run inference on extremely small models.
Microsoft's TinyStories~\cite{eldan2023tinystories} showed that models with as
few as 2.5M parameters can produce grammatically correct and contextually
consistent text when trained on a carefully curated dataset. Meanwhile,
the BitNet b1.58 architecture~\cite{ma2024bitnet} proved that ternary
weights (\{-1, 0, +1\}) can match full-precision performance at dramatically
reduced computational cost, and Zhu et al.'s Matmul-free
LLM~\cite{zhu2024matmulfree} eliminated matrix multiplication entirely
through scalable bitwise operations. On the embedded front,
MambaLite-Micro~\cite{mambalite2024} validated state space model (SSM)
inference on \espss{} with 83\% memory reduction over Transformer baselines.

Despite these advances, a critical question remains unanswered:
\emph{What is the minimum model capacity required to generate coherent
English text within the memory constraints of a commodity microcontroller?}
No existing work systematically establishes this lower bound or demonstrates
how architecture-specific hardware optimizations can push it further.

In this paper, we present \tinybit{}, a framework that addresses this gap
through two complementary contributions:

\begin{enumerate}[leftmargin=*,topsep=2pt,itemsep=1pt]
  \item \textbf{Language Coherence Lower Bound Analysis.} We train a family of
        GPT-2-style~\cite{radford2019language} models at four scales
        (1M, 3M, 8M, and 15M parameters) on the TinyStories dataset and
        systematically evaluate their language coherence under progressive
        quantization (FP32 $\rightarrow$ INT8 $\rightarrow$ INT4). We propose
        a multi-metric evaluation framework combining perplexity, MAUVE
        score~\cite{pillutla2021mauve}, and LLM-as-judge assessment to
        establish a Pareto frontier between model capacity, quantization level,
        and coherence quality.

  \item \textbf{Xtensa-Native Bitwise Inference Kernel.} We design a ternary
        General Matrix-Vector multiplication (\gemv{}) kernel that exploits the
        Xtensa LX7's SIMD extensions (EE instruction set) to perform
        inference without floating-point multiplication. By encoding ternary
        weights as packed 2-bit representations and decomposing \gemv{} into
        bitwise AND, XOR, and population count operations, we achieve
        significant speedup over conventional INT8 inference on the \espss{}.
\end{enumerate}

Together, these contributions enable end-to-end language model inference on the
\espss{} with a deployment footprint under 8\,MB, demonstrating coherent
English text generation on commodity microcontroller hardware for the first
time at meaningful model scales.

The remainder of this paper is organized as follows. Section~\ref{sec:related}
reviews related work in efficient LLMs and MCU deployment. Section~\ref{sec:problem}
formalizes the constrained inference problem. Section~\ref{sec:coherence}
presents our coherence lower bound experiments. Section~\ref{sec:kernel}
describes the Xtensa-native kernel design. Section~\ref{sec:evaluation}
reports end-to-end evaluation results. Section~\ref{sec:discussion} discusses
implications and limitations, and Section~\ref{sec:conclusion} concludes.
